\documentclass[10pt,a4paper]{amsart}
\usepackage[margin=19mm]{geometry}
\usepackage{amsmath,amssymb,mathtools}
\usepackage[scr=rsfs,cal=euler]{mathalfa}
\usepackage{xfrac}
\usepackage[unicode,hidelinks,bookmarks=false]{hyperref}
\newcommand{\ie}{i.e.\ }
\newcommand{\cf}{cf.\ }
\newcommand{\resp}{respectively\ }
\newcommand{\Subsection}{\subsection}
\providecommand{\nobreakdash}{\nobreak\hbox{-}}
\setlength{\emergencystretch}{2em}
\multlinegap0pt
\usepackage{bm}
\usepackage[shortlabels]{enumitem}
\usepackage{stackengine}
\usepackage{amssymb}
\newtheorem{lemma}{Lemma}
\title{Hadamard Hypercubes}
\author{Amin Bahmanian, Sho Suda}
\date{Source: arXiv:2605.16722v1 (CC BY 4.0)}
\begin{document}
\maketitle

\noindent\emph{Source and licence.} Hadamard Hypercubes. Version arXiv:2605.16722v1 (CC BY 4.0), distributed by arXiv under the Creative Commons Attribution 4.0 International licence, http://creativecommons.org/licenses/by/4.0/, as recorded in the arXiv licence metadata for this exact version and shown on the abstract page https://arxiv.org/abs/2605.16722v1. Full licence notice: \texttt{licenses/arxiv-2605.16722-CC-BY-4.0.txt}.\par\smallskip

\noindent\textit{Editorial scope.} Counted unit: the article's lemma lem:regularskewtwograph (source0.tex lines 572--578). The article's complete original proof of it is delivered with it (source0.tex lines 579--634, one \texttt{proof} environment of 56 lines). No mathematics is written, repaired or re-derived: the statement and the proof are the article's own lines, in the article's own notation, and no retained line is substituted or re-flowed.  No further source-local context is needed: every cross-reference of every retained line resolves inside the two delivered spans.  Nothing was omitted from any retained span, so the package carries no omission record.  No retained line contains a citation, so the wrapper carries no bibliography and no bibliography entry.  No retained line contains a figure environment, an image-inclusion command or drawing code, so no image is delivered and the package declares no asset dependency.  Editorial formatting only, in addition: an AMS \texttt{amsart} preamble that restates the article's own declarations and macros used by the retained lines, namely the environments \texttt{lemma} on separate counters, together with the standard packages those lines need.  The retained environments print in this document as Lemma 1 in order of appearance, whereas the article prints the same items with the numbering of its own full document.  The complete native source of the article as distributed in the arXiv e-print for this exact version is bundled unchanged as \texttt{sources/200742/source0.tex}.
\begin{lemma}\label{lem:regularskewtwograph} 
For a skew two-graph $(X,\nabla)$  and  $w\in X$,  the following are equivalent. 
    \begin{enumerate}
        \item [\textup{(i)}] The skew two-graph $(X,\nabla)$ is  regular. 
        \item [\textup{(ii)}] The descendent tournament $T_w$ is doubly regular. 
    \end{enumerate}
\end{lemma}

\begin{proof}
    Set $n=|X|$. 
   First, let us assume that (i)  holds.     Since $(X,\nabla)$ is regular, 
    the size of the following set 
    \[
    \{z\in X \mid (wyz)\in\nabla\}
    \]
    for $y\neq x$ equals the degree, $(n-2)/2$ where $n=|X|$, which is also equal to the number of edges $(y,z)$ in $T_w$. 
    Therefore $T_w$ is a regular tournament for any $w$. 
    Now, we fix two distinct vertices $y,z$ in $T_w$ and we may assume that $(y,z)\in E_w$. Let
    \begin{align*}
        N_w(x)&=\{u\in X\setminus\{w\} \mid (x,u)\in E_w\},\\
        \overline{N}_w(x)&=\{u\in X\setminus\{w\} \mid (x,u)\not\in E_w\}, 
    \end{align*}
    for $x\in X\setminus\{w\}$. 
    Since $\mathcal{T}_w$ is regular, 
    \begin{align*}
        |N_w(x)|=|\overline{N}_w(x)|=\frac{n-2}{2}    
    \end{align*}
    for any $x\in X\setminus\{w\}$. 
    Since $|N_w(y)|=1+|N_w(y)\cap N_w(z)|+|N_w(y)\cap \overline{N}(z)|$ and $|\overline{N}_w(y)|=|\overline{N}_w(y)\cap N(z)|+|\overline{N}_w(y)\cap \overline{N}_w(z)|$,  we have 
    \begin{align}\label{eq:drt1}
        1+|N_w(y)\cap N_w(z)|+|N_w(y)\cap \overline{N}_w(z)|=|\overline{N}_w(y)\cap N_w(z)|+|\overline{N}_w(y)\cap \overline{N}_w(z)|=\frac{n-2}{2}.  
    \end{align}
    Similarly, by considering $N_w(z)$, 
    \begin{align}\label{eq:drt11}
        |N_w(y)\cap N_w(z)|+|\overline{N}_w(y)\cap N_w(z)|=\frac{n-2}{2}.  
    \end{align}
    

    We now consider another descendent tournament $T_y$. Since $T_y$ is obtained from the tournament $T_w \cup\{w\}$, which is a tournament obtained from $T_w$ by adding a dominating vertex $w$, by switching with respect to $\{y\}\cup N_w(y)$ and deleting the resulting dominating vertex $y$.  
    We now apply the same argument to this tournament $T_y$ as above. 
    Since $T_y$ is regular, 
    \begin{align*}
        |N_y(z)|=\frac{n-2}{2}.     
    \end{align*}
    On the other hand, it holds that 
    \begin{align*}
        N_y(z)&=(\overline{N}_w(y)\cap N_w(z))\cup(N_w(y)\cap \overline{N}_w(z)).  
    \end{align*}
    Therefore we have 
    \begin{align}\label{eq:drt2}
        |\overline{N}_w(y)\cap N_w(z)|+|N_w(y)\cap \overline{N}_w(z)|=\frac{n-2}{2}.  
    \end{align}
    Combining \eqref{eq:drt1}, \eqref{eq:drt11}, \eqref{eq:drt2} yields that 
    \begin{align*}
        |N_w(y)\cap N_w(z)|&=\frac{n-4}{4},\\
        |N_w(y)\cap \overline{N}(z)|&=\frac{n-4}{4},\\
        |\overline{N}_w(y)\cap N(z)|&=\frac{n}{4},\\
        |\overline{N}_w(y)\cap \overline{N}_w(z)|&=\frac{n-4}{4}.
    \end{align*}
    Therefore, $|N_w(y)\cap N_w(z)|$ does not depend on the choice of $y,z$, and (ii) holds. 
    
    (ii) $\Rightarrow$ (i): 
    Assume (ii) holds. Each tournament $T_w$ is regular. Therefore, the valency of the skew two-graph is constant, (i) holds.   
\end{proof}

\end{document}
